\documentclass[10pt]{article}
\usepackage[utf8]{inputenc}
\usepackage[T1]{fontenc}
\usepackage{amsmath}
\usepackage{amsfonts}
\usepackage{amssymb}
\usepackage[version=4]{mhchem}
\usepackage{stmaryrd}

\begin{document}
странство $X$ наз. $p$-пространством, если оно имеет 0 . в своем Стоуна - Чеха бикомпактном расширении или Уолмена бикомпактиом расиирении. Каждое полное в смысле Јеха пространство является $p$-пространством. Каждое $p$-пространство имеет точечно-счетный тип. В $p$-пространстве справедливы аддиционная теорема для веса, и сетевой вес совпадает с весом. Паракомпактные $p$-пространства - это в точности совершенные прообразы метрич. пространств. Паракомпактные $p$ пространства с точечно-счетной базой метризуемы, равно как метризуемы и пространства этого вида с диагональю $G_{\delta}$. Совершенный образ и совершенный прообраз паракомпактного $p$-пространства - также паракомпактные $p$-пространства. B. И. Погомарев.

ОПОРНАЯ ГИПЕРПЛОСКОСТЬ м Н о ж е с т в а $M$ в $n$-мер ном век т о р ном п р ос т р ан с тв е - $(n-1)$-мерная плоскость, к-рая содержит точки замыкания $M$ и оставляет $M$ в одном замкнутом пространстве. При $n=3$ О. Г. наз. о п о р н о й п Л ос костью, а при $n=2-0$ п о рн ой прямой.

Граничную точку множества $M$, через к-рую проходит хотя бы одна О.. ., наз. о II о р н о їі т о ч к о ї̈ $M$. У выпуклого множества $M$ все его граничные точки опорные. Последнее свойство Архимед использовал как определение выпуклости $M$. Граничные точки выпуклого множества $M$, через к-рые проходит единственная О.г., наз. г л а д к и м и.

В общих векторных пространствах, где гиперплоскость определяется как область постоянства значений линейного функционала, также вводится понятие О. г. как гишерплоскости, әкстремальной по значению этого функционала среди гиперплоскостей, оставляющих $M$ в одном полупространстве.

ОПОРНАЯ ФУНКЦИЯ, о п о р ны й ф ун к ц и ов а л, множества $A$, лежащего в векторном пространстве $X$,- функция $s A$, задаваемая в находящемся с ним в двойственности векторном пространстве $Y$ соотношением

\[
(s \boldsymbol{A})(y)=\sup _{y \in A}\langle x, y\rangle
\]

Напр., О. Ф. единичного гшара в нормированном пространстве, рассматриваемом в двойственности со своим сопряженным пространством, - это норма в последнем.

О. ф. всегда выпуклая, замкнутая и положительно однородная (первой степени). Оператор $s: A \rightarrow s A$ взаимно однозначно отображает совокупность выпуклых замкнутых множеств в $X$ на совокупность выпуклых замкнутых однородных функций, обратный оператор не что иное, как субдифберенциал (в нуле) опорной функции. Именно, если $A$ - выпуклое замкнутое подмножество в $X$, то $\partial(s A)=A$, и если $p-$ выпуклая замкнутая однородная функция на $Y$, то $s(\partial p(0))=p$. Эти два соотношения (я вляющиеся следствием теоремы Фенхеля - Моро, см. Сопряженная функция) и выражают двойственность между замкнутыми выпуклыми множествами и выпуклыми замкнутыми однородными функциями.

Примеры соотношений, связывающих оператор $s$ c алгебраическими и теоретико-множественными операциями:

\[
\begin{aligned}
& s(\lambda C)=\lambda s C, \lambda>0 ; s\left(A_{1}+A_{2}\right)=s A_{1}+s A_{2} ; \\
& s\left(\operatorname{conv}\left(A_{1} \cup A_{2}\right)\right)(x)=\max \left(s A_{1}(x), s A_{2}(x)\right) .
\end{aligned}
\]

Лит. [1] Р о қ а ф е лі л а p Р., Выпуклый анализ, пер. c arT.I., M., 1973; 12] M i n k o w s k i H., Geometrie der Zahlen, Lpz.- B., 1911; L4] F e n c $\mathrm{h}$ e $1 \mathrm{~W}$., "Canad. J. Math.», 1949, v. 1, p. 73-77; [5] e r o \% e, Convex cones, sets and funktions, Princeton, 1953; [6] H ö r m a n d e r L., "Ark. för Mat.", 1955, bd 3, p. $181-86$.

B. М. Тихомиров.

ОПРЕДЕЛЕННЫИ ИНТЕГРАЛ - сМ. Интеграब. ОПРЕДЕЛИТЕЛЬ, д е т е р м и в а н т, квадратной матрицы $A=\left\|a_{i j}\right\|$ порядка $n$ над ассоциативно-коммутативным кольцом $K$ с единицеї 1 - элемент кольца $K$, равный сумме всех членов вида

\[
(-1)^{t} a_{1 i_{1}} \ldots a_{n i_{n}}
\]

где $i_{1}, \ldots, i_{n}$ - перестановка чисел $1, \ldots, n$, а $t-$ число инверсий перестановки $i_{1}, \ldots, i_{n}$. О. матрицы

\[
A=\left\|\begin{array}{ccc}
a_{11} & \ldots & a_{1 n} \\
\cdots & \cdot & \cdot \\
a_{n 1} & \ldots & a_{n n}
\end{array}\right\|
\]

обозначается

\[
\left|\begin{array}{ccc}
a_{11} & \ldots & a_{1 n} \\
\cdot & \cdot & \cdot \\
a_{n 1} & \ldots & a_{n n}
\end{array}\right| \text { шли } \operatorname{det} A .
\]

О. матрицы $A$ содержит $n$ ! членов; при $n=1 \operatorname{det} A=$ $=a_{11}$, при $n=2 \operatorname{det} A=a_{11} a_{22}-a_{21} a_{12}$. Наиболее важные для приложений случаи: $K$ - поле (в частности, числовое поле), $K$ - кольцо функций (в частности, колыцо многочленов), $K$ - кольцо целых чисел.

Всюду ниже $K$ - ассоциативно-коммутативное кольцо с $1, M_{n}(K)$ - совокупность всех квадратных матриц порядка $n$ над $K, E_{n}$ - единичная матрица над $K$. Пусть $A \in M_{n}(K)$, а $a_{1}, \ldots, a_{n}$ - строки матрицы $A$ (все далее изложенное справедливо и для столбцов матрицы $A)$. О. матрицы $A$ удобно рассматривать как функцию от ее строк:

\[
\operatorname{det} A=D\left(a_{1}, \ldots, a_{n}\right) \text {. }
\]

Отображение

\[
d: M_{n} \longrightarrow K(A \longmapsto \operatorname{det} A)
\]

подчинено следующим трем условиям:

\begin{enumerate}
  \item $d(A)$ - линейная функция любой строки матрицы $A$ :
\end{enumerate}

\[
\begin{gathered}
D\left(a_{1}, \ldots, \lambda a_{i}+\mu b_{i}, \ldots, a_{n}\right)= \\
=\lambda D\left(a_{1}, \ldots, a_{i}, \ldots, a_{n}\right)+\mu D\left(a_{1}, \ldots, b_{i}, \ldots, a_{n}\right),
\end{gathered}
\]

где $\lambda, \mu \in K$

\begin{enumerate}
  \setcounter{enumi}{1}
  \item если матрица $B$ получена из $A$ заменой строки $a_{i}$ строкой $a_{i}+a_{j}, i \neq j$, то $d(A)=d(B)$;

  \item $d\left(E_{n}\right)=1$.

\end{enumerate}

Условия 1) - 3) однозначно определяют отображение $d$, т. е. если отображение $h: M_{n}(K) \rightarrow K$ удовлетворяет условиям 1) -3 ), то $h(A)=\operatorname{det} A$. Таким образом получается аксиоматич. построение теории $\mathrm{O}$.

Пусть отображение $f: M_{n}(K) \rightarrow K$ удовлетворяет условию:

$1_{\text {a) если }}^{B}$ получается из матрицы $A$ умножением одноӥ строки на $\lambda \in K$, то $f(B)=\lambda f(A)$. Очевидно,, 1$\left.) \Rightarrow 1_{\mathrm{a}}\right)$. В случае, когда $K$ - поле, совокупность условий 1) - 3) оказывается равносильной условиям $1_{a}$ ), 2), 3).

О. диагональной матрицы равен произведению ее диагональных элементов. Отсюда вытекает сюръективность отображения $d: M_{n}(K) \rightarrow K$. О. треугольной матрицы также равен произведению ее диагональных элементов. Для матрицы $A=\left\|\begin{array}{ll}B & 0 \\ D & C\end{array}\right\|, \quad$ где $B \quad$ и $\quad C-$ квадратные матрицы,

\[
\operatorname{det} A=\operatorname{det} B \operatorname{det} C .
\]

Из свойств перестановок вытекает, что $\operatorname{det} A^{T}=\operatorname{det} A$, где ${ }^{T}$ - знак транспонирования. Если матрица $A$ имеет две одинаковые строки, то ее определитель равен 0 ; если поменять местами две строки матрицы $A$, то ее $O$. измевит знак;

$D\left(a_{1}, \ldots, a_{i}+\lambda a_{j}, \ldots, a_{n}\right)=D\left(a_{1}, \ldots, a_{i}, \ldots, a_{n}\right)$ при $i \neq j, \lambda \in K$; для $A$ и $B$ из $M_{n}(K)$ $\operatorname{det} A B=\operatorname{det} A \operatorname{det} B$


\end{document}